The \peripheral{SARADC} is a 10-bit successive-approximation register (SAR) analog-to-digital converter based on a capacitive DAC (CAPDAC) architecture, optimized for compact layout and low energy consumption. A single trigger initiates one conversion cycle; continuous conversion mode (\texttt{CONTMEAS}) allows repeated conversions without further software intervention.

\subsubsection*{Resolution}

The ADC produces a 10-bit unsigned result, available in bits 9:0 of the \texttt{SARADC\_DATA} register. Bits 31:10 read as zero.

\subsubsection*{Conversion Timing}

The conversion cycle is synchronized to the peripheral clock and consists of three sequential phases: a clear phase, a sample phase, and a successive-approximation conversion phase. The duration of the sample phase is controlled by the \texttt{SAMPLESTEP} field, which sets the initial value of a 4-bit countdown counter that runs during the sample phase. A value of~0 gives the minimum sample time. Refer to the ADC characterization data for the relationship between sample step value and the resulting sample duration.

\subsubsection*{Interrupts and Flags}

Two status flags in \texttt{SARADC\_SR} indicate conversion state:

\begin{itemize}
    \item \textbf{DATAVALID} (SR bit~1): set by hardware when a new conversion result is available in \texttt{SARADC\_DATA}. Clear by writing~1 to this bit. If this flag is already set when a new conversion completes, the overflow flag \texttt{OVF} is also set.
    \item \textbf{OVF} (SR bit~2): set when a conversion completes while \texttt{DATAVALID} is still set, indicating that the previous result was not read in time. Clear by writing~1 to this bit.
\end{itemize}

An interrupt is asserted when \texttt{DATAVALID} is set and the data-valid interrupt enable bit \texttt{DATAIE} is~1.

\subsubsection*{Operation}

\begin{enumerate}
    \item Set \texttt{SAMPLESTEP} to the desired sample phase duration (0 for minimum).
    \item Optionally set \texttt{CONTMEAS} to~1 for continuous conversion mode, and/or set \texttt{DATAIE} to~1 to enable the data-valid interrupt.
    \item Enable the ADC by setting \texttt{EN} to~1 in \texttt{SARADC\_CR}. This starts the first conversion.
    \item Wait for conversion to complete by polling \texttt{BUSY} until it reads~0, or wait for the IRQ if interrupts are enabled.
    \item Read the 10-bit result from bits 9:0 of \texttt{SARADC\_DATA}.
    \item Clear \texttt{DATAVALID} by writing~1 to SR bit~1. In continuous mode, the next conversion begins automatically after the previous one completes.
\end{enumerate}

\subsubsection*{Digital Test Ports}

The \texttt{SARADC\_TPR} (Test Port Register) provides a hardware debugging aid that routes internal ADC signals to two external digital test pins, \texttt{DTP0} and \texttt{DTP1}. These pins can be probed with a logic analyser or oscilloscope to observe the internal state of the peripheral in real silicon, where internal nets are otherwise inaccessible. The test ports are only active when debug mode is enabled (\texttt{DEBUG}~=~1); both outputs are driven to~0 otherwise.

The 4-bit field \texttt{DTP0SEL} (bits 3:0 of \texttt{SARADC\_TPR}) selects which internal debug signal is driven onto \texttt{DTP0}; the 4-bit field \texttt{DTP1SEL} (bits 7:4) selects the signal for \texttt{DTP1}. Both fields use the same index space, so any two signals can be observed simultaneously. The available debug signal indices are listed in Table~\ref{t:saradc-dtp}.

\begin{table}[h]
\centering
\begin{tabularx}{\textwidth}{c l X}
\hline
\textbf{Index} & \textbf{Signal} & \textbf{Description} \\
\hline
0  & \texttt{SAMPLESTEP[0]} & Bit~0 (LSB) of the \texttt{SAMPLESTEP} field currently loaded from \texttt{SARADC\_CR}. \\
1  & \texttt{SAMPLESTEP[1]} & Bit~1 of \texttt{SAMPLESTEP}. \\
2  & \texttt{SAMPLESTEP[2]} & Bit~2 of \texttt{SAMPLESTEP}. \\
3  & \texttt{SAMPLESTEP[3]} & Bit~3 (MSB) of \texttt{SAMPLESTEP}. \\
4  & \texttt{EN}            & ADC enable bit (\texttt{CR} bit~5). High when the ADC is enabled. \\
5  & \texttt{DATAIE}        & Data-valid interrupt enable bit (\texttt{CR} bit~7). \\
6  & \texttt{CONTMEAS}      & Continuous measurement mode enable bit (\texttt{CR} bit~8). \\
7  & \texttt{ADC\_ready}    & Registered copy of the ready strobe from the analog ADC core. Goes high on each cycle where new data is presented. \\
8  & \texttt{BUSY}          & Conversion busy flag (\texttt{SR} bit~0). High while the SAR conversion cycle is in progress. \\
9  & \texttt{OVF}           & Overflow interrupt flag (\texttt{SR} bit~2). High when a conversion completed while \texttt{DATAVALID} was already set. \\
10 & \texttt{DATAVALID}     & Data-valid flag (\texttt{SR} bit~1). High when a conversion result is ready to be read from \texttt{SARADC\_DATA}. \\
11 & \texttt{clr\_ovf}      & Internal clear-OVF pulse. Briefly high when software writes~1 to \texttt{SR} bit~2 to clear the overflow flag. \\
12 & \texttt{clr\_valid}    & Internal clear-DATAVALID pulse. Briefly high when software writes~1 to \texttt{SR} bit~1 to clear the data-valid flag. \\
13--15 & ---                & Tied low. Not connected to any internal signal. \\
\hline
\end{tabularx}
\caption{SARADC digital test port debug signal index map.}
\label{t:saradc-dtp}
\end{table}
